\rhead{ML-EL3: Generative AI \& Large Language Models}
\phantomsection 
\section*{Generative AI \& Large Language Models (ML-EL3)}
\label{elec:ML_3}

\noindent \small{\hyperref[main:scheme]{$\leftarrow$ Back to Scheme}} \vspace{0.5cm}

\noindent \textbf{Credits:} 4 (3-0-2)

\subsection*{Theory}
\noindent\textbf{Unit 1} \\
\textit{Generative Modeling Foundations:} Generative vs. discriminative models, Maximum likelihood estimation for density estimation, Autoregressive models and sequential generation, Latent variable models, Evaluation metrics (log-likelihood, FID, IS scores), Mode collapse and posterior collapse problems, Training instabilities and divergence measures.

\vspace{0.3cm}

\noindent\textbf{Unit 2} \\
\textit{Variational Autoencoders and Diffusion Models:} VAE architecture (encoder-decoder, KL divergence), Reparameterization trick and $\beta$-VAE variants, Hierarchical VAEs and VQ-VAE, Denoising Diffusion Probabilistic Models (DDPM), Forward/reverse diffusion process, Denoising score matching, Classifier-free guidance, Latent diffusion models (Stable Diffusion).

\vspace{0.3cm}

\noindent\textbf{Unit 3} \\
\textit{Large Language Model Architectures:} Decoder-only transformers (GPT series evolution), Rotary Position Embeddings (RoPE), Grouped-Query Attention (GQA), Mixture of Experts (MoE) scaling, Context window extension techniques (RoPE scaling, ALiBi), FlashAttention and memory-efficient attention, Speculative decoding and key-value caching optimizations.

\vspace{0.3cm}

\noindent\textbf{Unit 4} \\
\textit{Pretraining, Alignment, and Scaling Laws:} Self-supervised pretraining objectives (causal LM, masked LM, prefix LM), Chinchilla scaling laws and compute-optimal training, Instruction tuning and dataset curation, Reinforcement Learning from Human Feedback (RLHF), Proximal Policy Optimization (PPO) for alignment, Direct Preference Optimization (DPO), Constitutional AI and self-improvement.

\vspace{0.3cm}

\noindent\textbf{Unit 5} \\
\textit{Generative AI Applications and Deployment:} Retrieval-Augmented Generation (RAG) architectures, Agentic workflows (ReAct, Toolformer), Multimodal generation (CLIP, DALL-E, Flamingo), Text-to-image/video/speech generation, Guardrails and content moderation, Model evaluation (human preference, Elo rankings), Production deployment (vLLM, TGI, quantization, distillation), Ethical considerations and safety measures.

\newpage

\subsection*{Practical}
\noindent\textbf{Lab 1: VAE Implementation} \\
Focuses on complete VAE training, latent space visualization, and $\beta$-VAE experimentation.

\vspace{0.2cm}

\noindent\textbf{Lab 2: Diffusion Model Training} \\
Focuses on DDPM implementation, sampling visualization, and classifier guidance.

\vspace{0.2cm}

\noindent\textbf{Lab 3: Transformer from Scratch} \\
Focuses on GPT-like decoder implementation with RoPE and FlashAttention.

\vspace{0.3cm}

\noindent\textbf{Lab 4: Fine-tuning Small Language Models} \\
Focuses on instruction tuning datasets and LoRA parameter-efficient fine-tuning.

\vspace{0.2cm}

\noindent\textbf{Lab 5: RLHF with PPO/DPO} \\
Focuses on preference dataset creation and alignment training.

\vspace{0.2cm}

\noindent\textbf{Lab 6: RAG Pipeline Development} \\
Focuses on dense retrieval, reranking, and generation with knowledge augmentation.

\vspace{0.2cm}

\noindent\textbf{Lab 7: Multimodal Generation} \\
Focuses on CLIP-guided diffusion and text-to-image synthesis.

\vspace{0.2cm}

\noindent\textbf{Lab 8: LLM Agent Development} \\
Focuses on ReAct agents with tool calling and reasoning capabilities.

\vspace{0.2cm}

\noindent\textbf{Lab 9: Production LLM Serving} \\
Focuses on vLLM/TGI deployment, quantization, and batching optimization.

\vspace{0.2cm}

\noindent\textbf{Lab 10: Generative AI Application} \\
Focuses on complete production generative system with evaluation and safety measures.

\vspace{0.2cm}

\newpage
\rhead{Curriculum Schema}
